\documentclass[a4paper,11pt]{article}

\usepackage{revy}
\usepackage[utf8]{inputenc}
\usepackage[T1]{fontenc}
\usepackage[danish]{babel}
\usepackage{amsmath,amssymb}

\revyname{MatRevy}
\revyyear{2026}
\version{2026-10-09}
\eta{3.5 minutter}

\title{Beskidt Matematik}
\author{Markus '20}

\begin{document}
\maketitle

\begin{roles}
\role{S1}[Skuespiller] Skuespil
\role{S2}[Skuespiller] lille person
\role{OE}[Skuespiller] Onde Ernst
\role{N1}[Ninja] Ninja, stærk person
\role{N2}[Ninja] Ninja, stærk person
\end{roles}

\begin{sketch}
\scene{S1 og S2 er almindelige studerende. Onde Ernst er klædt som en tegneserie skurk ala Dracula. OE har magiske krafter og kan få matematiske ting til at flyve, som ninjaerne bærer. Evt. græskebogstaver eller udtryk}

\says{S1} Øj matintro er bare så spændende, men altså det med studieretningerne forvirerer mig.
\says{S2} Hvad forvirer dig?
\says{S1}Jamen der er gymnasiel, statistik og så ren matematik?
\says{S2} Det giver da god mening?
\says{S1} Altså jaja, men nej.. altså ren matematik, hvorfor er det så rent? Findes der beskidt matematik???
\scene{*chok lyd fra band og S2*}
\scene{Mørkt bag tæppe går til siden, OE træder frem…}
\says{OE} hmhmmh var der nogen der sagde… beskidt matematik?
\says{S1 og S2} øøhhh
\says{OE} var det dig!?\act{Peger på S1}
\says{S1}\act{Peger på S2}
\says{OE} \act{STORMER mod S2 tager fat i kraven og ved hjælp af to ninjaer bliver S2 løftet op i sin krave.}
\says{S2}\act{Fuld af rædsel i ansigtet.}
\says{S2}[Stammende] Hvad nej? Hvem er du.
\says{OE} hvem jer er?? \act{slipper S2 som bliver hængende i luften}
\says{OE} JEG ER ONDE ERNST!
\scene{*TORDEN SKRALD*}
\says{OE}[Griber fat i S2 igen] men det egentlige spørgsmål er ikke hvad beskidt matematik er. Det er om i er værdige til at blive åbnet op til en ny verden af overtællelige dimensioner.
\says{OE} \act{styrer S2 ned (som the force)}
\scene{Ninjaer sætter S2 langsomt ned}
\says{S1 og S2} \act{ kigger på hinanden i frygt}
\says{S1-2} Det ved vi ikke om vi er
\says{OE} Nej selvfølgeligt ved I ikke det. Det er noget jeg vurderer.
\says{OE} S1 hvornår har du sidst gjort rent i snullet?
\says{S1} uhhh... aldrig?
\says{OE} hmm goodt… godt… S2! Hvordan vil du tjekke om denne funktion er vel defineret? \act{Hiver en funktion frem?}
\says{S2} øhhh, det tror jeg faktisk slet ikke at jeg ville gøre.
\says{OE} KORREKT!
\says{OE} Nuvel, i viser lovende potientiale… Lad mig uddybe. Beskidt matematik, er alt det ufine. En stolt tradition udbygget af en lang række dygtige matematikkere tit overset af deres mere skole-rene tvillinger. Vi taler om Karl friedrich snavs, Wallodi Weibullshit og John von Nullermann. De har lavet det hårde arbejde, imens teoretikerne sidder i deres elfenbenstårn ude af trit med hvordan vi andre får verden til at dreje rundt.
\says{OE} \act{Står triumferende med armene til siden}
\says{S1} Vil det sige at Gauss og de andre har haft en ond tvillinge bror?
\says{S2} Ja findes der også en Kolmo-go-røv?
\says{OE} Var din mund! Dette er ikke noget at spøge med. \act{Peger på S2}
\says{N1} {dækker S2’s mund}
\says{OE} \act{Går truende mod S2}
\says{S1}[Prøver at lede OE væk fra S2] Hey, nå men hvad lærer man så i beskidt matematik.
\says{OE}\act{grynter} I første uge vil i lære at dividere med 0
\says{S2} \act{Gisp}
\says{OE}[henkastende] Ja ja trivielt, men bare se her
\scene{OE fører dem rundt i lokalet}
\says{OE} For at få dette bevis til at gå op skal vi bruge differentialerne som brøker
\says{S1} Men det kan man da ikke?
\scene{OE fører dem videre}
\says{OE} Her er udvalgs aksiomet ikke et valg! Det er tvang!
\says{S2} Men er det ikke farligt!?
\says{OE} Hahah… jo meget farligt, for de uværdige altså. Men i de rette hænder og med min vejledning kan du opnå uanede resultater! Selv Gödel sætninger vil kunne bevises!
\scene{OE bliver mere og mere ophidset}
\says{OE} I mit rige er pi og Eulers tal 3
\says{OE} $0^0$ er 2
\says{OE} Ikke bare beviserne er efterladt som en opgave til læseren, hele sætninger, NEJ HELE BØGER ER EFTERLADT SOM EN OPGAVE TIL LÆSEREN!
\scene{OE pruster og stønner af overanstrengelse}
\says{S1} Men vent, er det her ikke bare fysik?
\scene{Lys ned}
\end{sketch}
\end{document}
